\section{A uniform flow bound and finite-state closure}\label{sec:flow}

A total-degree estimate would only bound the sum of the exponents by a
quantity proportional to $n$. The finite-state construction instead
requires a bound at each position. The one-dimensional geometry of the
cyclic index set supplies that bound.

\begin{lemma}[Period-independent local flow bound]\label{lem:uniform-flow}
Let $n\ge1$, $d\ge2$, and $\mu\in\Z$. Suppose nonnegative integers
$L_i,R_i$, indexed cyclically, satisfy
\begin{equation}\label{eq:abstract-flow}
 d(L_i+R_i)\le\mu+R_{i-1}+L_{i+1}
 \qquad(0\le i<n).
\end{equation}
If $\mu<0$, no such configuration exists. If $\mu=0$, every exponent
is zero. If $\mu>0$, then for every $i$,
\begin{equation}\label{eq:uniform-bound}
 L_i+R_i\le\frac{\mu(d+1)}{(d-1)^2}.
\end{equation}
\end{lemma}

\begin{proof}
Set $e_i=L_i+R_i$. Summing \eqref{eq:abstract-flow} over the cyclic
index set gives
\begin{equation}\label{eq:total-flow}
 (d-1)\sum_i e_i\le n\mu.
\end{equation}
This is impossible for $\mu<0$. For $\mu=0$, nonnegativity forces
all $e_i$, and hence all $L_i,R_i$, to vanish.

Assume now that $\mu>0$ and fix one configuration. Construct an
$n\times n$ nonnegative matrix $P$ column by column. If $e_j>0$, put
$L_j/e_j$ in row $j-1$ of column $j$ and $R_j/e_j$ in row $j+1$
of that column. Add the two contributions if the rows coincide. If
$e_j=0$, put $1$ in row $j+1$ of column $j$. All unspecified entries
are zero. Each column has sum one, and the zero-flow convention does
not affect $P\ee$. Specifically,
\begin{equation}\label{eq:redistribution}
 (P\ee)_i=R_{i-1}+L_{i+1}.
\end{equation}
The inequalities become $d\ee\le\mu\one+P\ee$, componentwise.
Because $P$ is nonnegative, iterating this inequality $K$ times gives
\begin{equation}\label{eq:iterated-flow}
 \ee\le\frac{\mu}{d}\sum_{k=0}^{K-1}d^{-k}P^k\one
         +d^{-K}P^K\ee.
\end{equation}

Every power of $P$ is column-stochastic. Thus each of its entries lies
between zero and one, and
\begin{equation}\label{eq:tail-flow}
 \|P^K\ee\|_\infty\le\|\ee\|_1.
\end{equation}
Here $n$ and the configuration are fixed. Since $d>1$, the last term
in \eqref{eq:iterated-flow} tends to zero as $K\to\infty$. There is
no need to bound $\|\ee\|_1$ uniformly in $n$ at this stage.

To control the first term, fix a row $i$. The support of $P$ allows
only steps between adjacent cyclic positions. If $(P^k)_{ij}\ne0$,
then the cyclic distance between $i$ and $j$ is at most $k$. There are
at most $\min(n,2k+1)$ such starting columns $j$. Using the entrywise
bound from column-stochasticity gives the row-sum estimate
\begin{equation}\label{eq:row-sum}
 (P^k\one)_i=\sum_j(P^k)_{ij}\le2k+1.
\end{equation}
In particular, the proof does not mistake a column-stochastic matrix
for a row-stochastic one. The same estimate is valid for $n=1,2$,
where the adjacent positions may coincide.

Take $K\to\infty$ in \eqref{eq:iterated-flow} and use
\eqref{eq:row-sum}. Then
\begin{align}
 e_i&\le\frac{\mu}{d}\sum_{k\ge0}\frac{2k+1}{d^k}\notag\\
    &=\frac{\mu}{d}
      \left(\frac{d}{d-1}+\frac{2d}{(d-1)^2}\right)
      =\frac{\mu(d+1)}{(d-1)^2},\label{eq:geometric-flow-bound}
\end{align}
which proves \eqref{eq:uniform-bound}.
\end{proof}

The auxiliary matrix $P$ depends on the exponent configuration. It is
only a proof device, not a Markov model of the H\'enon dynamics.
Its entries are nonnegative even when $p,a,q$ are complex. The
quadratic case $d=2$ is covered without a denominator $d-2$.

\begin{proof}[Proof of Theorem~\ref{thm:main}]
Lemma~\ref{lem:normal-form} gives the finite scheme and defines the
residue for all parameters. By Lemma~\ref{lem:local-expansion}, each
nonzero product in \eqref{eq:configuration-sum} satisfies
\eqref{eq:abstract-flow} with $\mu=m-d+1$.
Lemma~\ref{lem:uniform-flow} shows that either all products vanish,
only the zero configuration contributes, or
\begin{equation}\label{eq:box-cutoff}
 0\le L_i,R_i\le L_i+R_i\le B
\end{equation}
for every $i$. The integer part in \eqref{eq:state-bound} is justified
because $L_i+R_i$ is an integer. Hence every nonzero term of the
convergent residue sum lies in a single finite box independent of $n$.

Associate to a cyclic configuration the edge states
\begin{equation}\label{eq:edge-bijection}
 E_i=(R_{i-1},L_i),\qquad E_{i+1}=(R_i,L_{i+1}).
\end{equation}
If $E_i=(r,\ell)$ and $E_{i+1}=(s,v)$, the $i$th local factor in
\eqref{eq:configuration-sum} is
\[
 \binom{\ell+s}{\ell}a^\ell c_{r+v,\ell+s}
 =W_{E_i,E_{i+1}}.
\]
Conversely, any labelled closed sequence of edge states determines
the $L_i$ as its second coordinates and the $R_i$ as the first
coordinates of the following states. This is inverse to
\eqref{eq:edge-bijection}, including for $n=1,2$. There is no quotient
by cyclic rotation. Consequently
\begin{align}
 \rho_n
 &=\sum_{E_0,\ldots,E_{n-1}\in\mathcal E}
       W_{E_0,E_1}W_{E_1,E_2}\cdots W_{E_{n-1},E_0}\notag\\
 &=\Tr(W^n).\label{eq:trace-closure}
\end{align}
Additional states or configurations inside the box with zero product
do not affect this equality. When $m-d+1\le0$, $B=0$ and the same
identity follows from the vanishing or the unique zero configuration
already established.

For the coefficient ring, set $T=1/X$. The reversed monic polynomial is
\[
 \widetilde p(T)=T^dp(1/T)
    =1+p_{d-1}T+\cdots+p_0T^d.
\]
Every coefficient of $\widetilde p(T)^{-(e+1)}$ is an
integer-coefficient polynomial in the $p_j$. Indeed, expand
$(1+U)^{-(e+1)}$ in powers of
$U=p_{d-1}T+\cdots+p_0T^d$, whose constant term is zero. A given
$T$-coefficient receives only finitely many terms, each with an
integer binomial coefficient. Formula~\eqref{eq:local-coefficient}
therefore gives an integer-coefficient polynomial in the $p_j,q_j$.
The extra binomial factor and $a^\ell$ in \eqref{eq:edge-matrix}
prove \eqref{eq:coefficient-ring}.

Definition~\ref{def:residue-trace} now gives
$\tau_n=-\Tr(W^n)$. The finite-dimensional identity
\begin{equation}\label{eq:logdet}
 -\log\det(I-zW)=\sum_{n\ge1}\frac{\Tr(W^n)}{n}z^n
\end{equation}
holds as a formal series. Substitution into
\eqref{eq:generating-function} yields the reciprocal determinant in
\eqref{eq:main}. Both the residue coefficients and the entries of
$W$ are polynomials, so the coefficient identities specialize at any
parameter, including a drop in the actual degree of $q$.
Finally $\det(I-zW)$ has constant term one, hence its reciprocal is
analytic near $z=0$.
\end{proof}

\begin{remark}[The low-degree threshold]\label{rem:threshold}
For $m<d-1$, the negative-surplus case gives $\rho_n=\tau_n=0$ for
all $n$. For $m=d-1$, only $L_i=R_i=0$ contributes and
$c_{0,0}=q_{d-1}$. Thus
\[
 \rho_n=q_{d-1}^{n},\qquad
 \tau_n=-q_{d-1}^{n},\qquad
 D_{p,a,q}(z)=\frac{1}{1-q_{d-1}z}.
\]
The subcritical vanishing is the classical Euler--Jacobi degree
phenomenon \citep[Corollary~1.18 of the author preprint]{cattani1996residues}.
It is not the additional uniform-flow assertion needed
for larger $m$.
\end{remark}
